\chapter{Process Discovery Simulation Methodology}
\label{chp:process_discovery_methodology}

\textit{erklärt, wie aus den Daten ein ausführbares Model wird und wie es arbeitet? | Modellkonstruktion und Ausführungsmechanismen erklären}

\section{Process Discovery}
Process Discovery Algorithms can be grouped into a family like structure. Each Family follows a methodology of how from observed sequnces, a modelstructure is constructed. The categories are overlapping in their calculation methods. 
\cite{DiscoveryBenchmark}.

    \textit{
    \begin{enumerate}
        \item Candidates and Configuration of the different Discovery Algorithms
        Inductive Miner / IMf

        Heuristics Miner
        
        ILP Miner
        
        Split Miner

        \item Procedures for Rule checks and ggf. Model adaptions
    \end{enumerate}}

\section{Prosits Simulation Methodology}
    \begin{enumerate}
        \item 
        \item Ankünfte, Routing, Ressourcenzuordnung, Kalender und Kapazitäten
        \item gelernte Zeitmodelle und Workload-Features
        \item Reihenfolge der Entscheidungen während der Simulation
        \item zusätzliche Diagnoseausgaben und Verifikationsprüfungen
    \end{enumerate}

\textit{Hier gehört ein nachvollziehbares Ausführungsbeispiel hinein: Ein Truck wird aktiviert, benötigt einen bestimmten Stop, erhält eine Ressource und durchläuft die einzelnen Zeitkomponenten. Jede Komponente muss zur Implementierung passen. So beantwortest du die Kritik an der bisher zu abstrakten technischen Darstellung.}